% UTF-8; compile directly with XeLaTeX, or input after loading xepersian.
% Faithful transcription; clinical statements have not been updated or corrected.
% Source images are preserved at neurology/neurology/images/pic01.jpg ... pic17.jpg.
\ifdefined\documentclass
  \ifdefined\XeTeXversion\else\errmessage{This file requires XeLaTeX}\fi
\fi
\ifdefined\settextfont
  \def\MeningitisStandalone{0}
\else
  \def\MeningitisStandalone{1}
  \documentclass[12pt,a4paper]{article}
  \usepackage[margin=2cm]{geometry}
  \usepackage{amsmath,amssymb,graphicx,xcolor,enumitem}
  \usepackage{xepersian}
  \IfFontExistsTF{Amiri}{\settextfont{Amiri}}{\settextfont{FreeSerif}}
  \setlatintextfont{FreeSerif}
\fi
\providecommand{\en}[1]{\lr{#1}}
\providecommand{\unclear}[1]{[خوانش نامطمئن: #1]}
\newcounter{menimage}
\newcommand{\menpage}[2]{\stepcounter{menimage}\clearpage\section*{صفحهٔ #1}\noindent\en{#2}\par\smallskip
\ifnum\value{menimage}<9
  \noindent\en{Subject: 100\quad Date: / /\quad YASHA}\par
\else
  \noindent\en{Date: / /\quad Sat. Sun. Mon. Tue. Thu. Wed. Fri.}\par
  \noindent\en{Subject:} مننژیت\par
\fi\medskip}
\newcommand{\menitem}[2]{\item[#1.] #2}
\ifnum\MeningitisStandalone=1 \begin{document}\fi
\menpage{۱}{pic01.jpg — 20261010\_053144594.jpg}
\section*{مننژیت — \en{chapter 100 -- Meningitis}}
\begin{description}
\menitem{۱}{التهاب لپتومننژ: باکتری، ویروس، قارچ (نادر)، انگل.}
\menitem{۲}{مننژیت آسپتیک: مننژیت ویروسی؛ کشت \en{CSF} منفی برای باکتری.}
\end{description}
\begin{itemize}
\item \en{Lyme disease; Syphilis; TB}.
\item عفونت‌های پارامننژیال: آبسهٔ مغزی، آبسهٔ اپی‌دورال، آمپیم ساب‌دورال.
\item داروها: \en{NSAID}، \en{IVIG}.
\item اختلالات خودایمنی.
\item بیماری کاوازاکی: واسکولیت عروق متوسط.
\end{itemize}
\begin{description}
\menitem{۳}{ارگانیسم‌های شایع قبل از واکسیناسیون: هموفیلوس آنفلوانزا، استرپ پنومونیه، نایسریا مننژیتیدیس.}
\end{description}

\menpage{۲}{pic02.jpg — 20261010\_053155765.jpg}
\begin{description}
\menitem{۴}{مننژیت استافیلوکوکی: \en{NICU} در نوزادان؛ ترومای نافذ به سر؛ جراحی نوروسرجری.}
\menitem{۵}{عوامل شایع مننژیت باکتریال:}
\end{description}
\begin{itemize}
\item نوزادان: \en{GBS}، \en{E. coli}، باسیل‌های گرم منفی روده‌ای.
\item ۱ تا ۳ ماه: استرپ پنومونیه، باسیل گرم منفی، \en{GBS}.
\item بالای ۳ ماه: استرپ پنومونیه، نایسریا مننژیتیدیس.
\end{itemize}
\begin{description}\menitem{۶}{\en{Partially treated meningitis}}
\end{description}
\begin{latin}\begin{itemize}
\item Bacterial meningitis $\to$ AB therapy before LP $\to$ Negative CSF culture.
\item Although other CSF findings suggestive of bacterial infection. $\to$ PCR of CSF.
\end{itemize}\end{latin}
\begin{description}\menitem{۷}{شایع‌ترین ویروس‌ها: انتروویروس، پارکوویروس؛ \en{HSV, EBV, CMV, HIV, LCMV}؛ \en{Mumps} در افرادی که واکسن نزده‌اند.}\end{description}

\menpage{۳}{pic03.jpg — 20261010\_053202786.jpg}
\begin{description}\menitem{۸}{\en{Less frequent infections causing meningitis}}\end{description}
\begin{latin}\begin{itemize}
\item Lyme disease (Borrelia burgdorferi).
\item Cat-scratch disease (Bartonella henselae).
\item TB.
\item Toxoplasma.
\item Fungi: Cryptococcus; Histoplasma.
\item Parasites: Angiostrongylus cantonensis; Naegleria fowleri.
\end{itemize}\end{latin}
\begin{description}
\menitem{۹}{آسپلنی آناتومیک یا عملکردی: \en{RF} مننژیت باکتریال.}
\menitem{۱۰}{\en{Crowding}: \en{RF} مننژیت باکتریال.}
\menitem{۱۱}{شکستگی قاعدهٔ جمجمه: \en{RF} مننژیت ناشی از استرپ پنومونیه.}
\menitem{۱۲}{مننژیت ممکن است به دنبال علائم دستگاه تنفس فوقانی رخ دهد.}
\menitem{۱۳}{علائم التهاب مننژ: سردرد، تحریک‌پذیری، تهوع و استفراغ؛ خواب‌آلودگی، فتوفوبی، سفتی گردن؛ معمولاً تب وجود دارد.}
\end{description}
\begin{latin}\begin{itemize}\item Kernig Sign; Brudzinski Sign: older than 12 months.\end{itemize}\end{latin}

\menpage{۴}{pic04.jpg — 20261010\_053212967.jpg}
\begin{description}
\menitem{۱۴}{در شیرخواران زیر ۱ سال علائم تحریک مننژ جزئی است.}
\end{description}
\begin{latin}\noindent Irritability, Depressed Mental Status; Poor feeding.\end{latin}
بقیهٔ علائم عصبی هم با شیوع کمتر ممکن است دیده شوند.
\begin{description}
\menitem{۱۵}{علائم افزایش \en{ICP} ممکن است دیده شود (در شیرخواران): \en{Bulging fontanelle}؛ \en{Diplopia, Ptosis}.}
\menitem{۱۶}{ادم پاپیل ناشایع است مگر: انسداد سینوس‌های وریدی، آمپیم ساب‌دورال، آبسهٔ مغزی.}
\menitem{۱۷}{شک به مننژیت باکتریال: \en{LP}. کنتراندیکاسیون: افزایش \en{ICP}؛ \en{Plt} زیر ۲۰٬۰۰۰؛ ناپایداری قلبی عروقی؛ آپنه.}
\menitem{۱۸}{\en{CSF Examination}: \en{WBC count and diff}؛ \en{Pr, Glucose level}؛ \en{Gram stain}؛ \en{Culture}؛ \en{if needed: PCR}.}
\end{description}

\menpage{۵}{pic05.jpg — 20261010\_053219366.jpg}
\begin{description}
\menitem{۱۹}{لکوسیتوز محیطی در مننژیت شایع است. \en{BC} ممکن است مثبت باشد.}
\menitem{۲۰}{\en{LP} باید قبل از درمان آنتی‌باکتریال انجام شود. اگر امکان \en{LP} نیست \en{AB} را به تعویق نمی‌اندازیم.}
\menitem{۲۱}{در مننژیت مننگوکوکی \en{CSF} حتی بعد از یک دوز \en{AB} و ظرف ۲–۱ ساعت استریل می‌شود؛ در استرپ پنومونیه چند ساعت.}
\menitem{۲۲}{\en{NL CSF}:}
\end{description}
\begin{latin}\begin{itemize}
\item P: $10$--$20\ \mathrm{cm\,H_2O}$.
\item $20$--$45\ \mathrm{mg/dl}$ pro.
\item $>50\%$ of serum glucose.
\item Up to $10$--$30\ \mathrm{WBC}/\mu\mathrm{L}$ is NL in neonates.
\item $\mathrm{WBC}<5$.
\end{itemize}\end{latin}
ارجحیت لنفوسیت.
\begin{description}\menitem{۲۳}{\en{CSF in Acute Bacterial Meningitis}:}\end{description}
\begin{latin}\begin{itemize}
\item P: $>25\ \mathrm{cm\,H_2O}$.
\item Pro: $100$--$500\ \mathrm{mg/dl}$.
\item Glucose: $<40\ \mathrm{mg/dl}$ or $<40\%$ serum glucose.
\item $\mathrm{WBC}>100$, PMN predominant.
\end{itemize}\end{latin}

\menpage{۶}{pic06.jpg — 20261010\_053229304.jpg}
\begin{description}
\menitem{۲۴}{هدف درمان در مننژیت باکتریال: استریل کردن \en{CSF}؛ حفظ پرفیوژن بافت مغز.}
\menitem{۲۵}{درمان امپایریک مننژیت باکتریال: نسل ۳ سفالوسپورین + ونکومایسین.}
\end{description}
\begin{itemize}
\item ونکومایسین: \en{S. pneumoniae}.
\item سفتازیدیم، سفپیم، سفتریاکسون: \en{N. meningitidis, Hi, E. coli}.
\item برای شیرخواران زیر ۲ ماه: آمپی‌سیلین هم می‌دهیم؛ \en{Listeria monocytogenes}.
\end{itemize}
\begin{description}
\menitem{۲۶}{طول درمان: ۵–۷ روز برای نایسریا؛ ۷–۱۰ روز برای \en{Hi}؛ ۱۰–۱۴ روز برای \en{S} پنومونیه.}
\menitem{۲۷}{داروی انتخابی:}
\end{description}
\begin{itemize}
\item نوزاد: \en{Cefta + Ampicillin $\pm$ Genta} یا \en{Ampicillin + Genta}.
\item ۱ ماه تا ۴ سال: \en{Ceftria + Vanco} یا \en{Cefepime + Vanco}.
\item ۵ تا ۱۲ سال: \en{Ceftria + Vanco} یا \en{Cefepime + Vanco}.
\end{itemize}

\menpage{۷}{pic07.jpg — 20261010\_053234720.jpg}
\begin{description}
\menitem{۲۸}{مننژیت ناشی از باسیل گرم منفی باید حداقل تا ۲۱ روز یا ۱۴ روز بعد از اولین کشت منفی \en{CSF} درمان شود. هر کدام که طولانی‌تر بود.}
\menitem{۲۹}{دگزامتازون را با دوز اول \en{AB} شروع می‌کنیم: \en{$\downarrow$ Hearing loss from Hi}.}
\menitem{۳۰}{\en{Supportive therapy for}: \en{Dehydration, shock, DIC, SIADH, Seizure, $\uparrow$ ICP, Apnea, Arrhythmias, Coma}.}
\menitem{۳۱}{از بین نرفتن تب طی درمان شایع است.}
\menitem{۳۲}{اگر حال عمومی بیمار ظرف ۴۸ ساعت بعد از درمان بهتر نشد \en{LP} مجدد انجام می‌دهیم؛ برای مننژیت‌های گرم منفی؛ بیماری که ایمنی‌اش سرکوب می‌شود.}
\menitem{۳۳}{اگر کودک علائم یا نشانه‌های فوکال نورولوژیک داشت و کشت مثبت پایدار: \en{CNS Imaging $\to$ CT with contrast / MRI}؛ افیوژن ساب‌دورال، آبسهٔ مغزی.}
\end{description}

\menpage{۸}{pic08.jpg — 20261010\_053246080.jpg}
\begin{description}
\menitem{۳۴}{هر بیماری مننژیت قبل از ترخیص: بررسی شنوایی.}
\menitem{۳۴}{عود مننژیت در ۳–۱۴ روز بعد از درمان شایع‌تر است. شک به یک ضعف سیستم ایمنی یا آناتومیک.}
\menitem{۳۵}{\en{Poor prognosis}: \en{young age, long duration, Seizure, Coma, shock, low CSF WBC in the presence of bacteria}.}
\menitem{۳۶}{افراد ناقل باید \en{AB} دریافت کنند: \en{Rifampine, Ciprofloxacin, Azithro, Ceftria}. نفوذ خیلی بالا [حاشیه: سل، بروسلوز].}
\menitem{۳۷}{کورتون نفوذ سفو به \en{CNS} کم می‌کند.}
\menitem{۳۸}{اگر کشت منفی بود: حداقل ۱۰ روز باید درمان شود.}
\menitem{۳۹}{اگر سفالوسپورین نگرفته باشد: باید بگیرد بعد از ترخیص.}
\menitem{۴۰}{اگر \en{SIADH}: $2/3$ مینیتننس مایع می‌دهیم.}
\end{description}

\menpage{۹}{pic09.jpg — 20261010\_053251011.jpg؛ شمارهٔ دست‌نویس: ۱۰}
\section*{مننژیت}
\begin{itemize}
\item ارگانیسم‌های شایع مننژیت باکتریال اکتسابی از جامعه: استرپ پنومونیه (۵۰\%)، نایسریا مننژیتیدیس (۲۵\%)، استرپ گروه \en{B} (۱۵\%)، لیستریا مونوسایتوژن (۱۰\%).
\item استرپ گروه \en{B} (آگالاکتیه): مننژیت در نوزادان، افراد بالای ۵۰ سال یا بیماری زمینه‌ای.
\item لیستریا مونوسیتوژن: مننژیت نوزادان، زنان باردار، سن بالای ۶۰، نقص ایمنی.
\item استرپ پنومونیه: شایع‌ترین عامل مننژیت در افراد بالای ۲۰ سال.
\end{itemize}
\subsection*{استرپ پنومونیه}
\begin{itemize}
\item دیپلوکوک گرم مثبت.
\item \en{RF}: پنومونی پنوموکوکی (مهم‌ترین \en{RF})؛ \en{OM} پنوموکوکی حاد و مزمن؛ الکلیسم؛ سینوزیت پنوموکوکی حاد و مزمن؛ \en{DM}؛ اسپلنکتومی؛ هایپوگاما؛ نقص کمپلمان؛ تروما به سر همراه با شکستگی \en{base} و نشت \en{CSF}.
\end{itemize}
\subsection*{نایسریا مننژیتیدیس}
\begin{itemize}
\item کوکسی گرم منفی هوازی.
\item واکسن آن: سروگروه \en{B} که عامل $1/3$ موارد است را شامل نمی‌شود.
\item واکسن سروگروه \en{B} مننگوکوک برای افراد ۱۶ تا ۲۳ سال توصیه می‌شود.
\item واکسن گروه \en{B} حاملین نازوفارنژیال را کاهش نمی‌دهد؛ در حالی که واکسن گروه \en{A, C, Y, W} حاملین نازوفارنژیال را کاهش می‌دهد.
\item عفونت ممکن است از کلونیزاسیون نازوفارنکس شروع شود.
\item بالینی کاراکتریستیک: پتشی و پورپورا.
\item \en{RF}: نقص در اجزای کمپلمان (\unclear{\en{C5--C9}}، پروپردین).
\item سیر بالینی: برق‌آسا؛ می‌تواند ظرف چند ساعت منجر به مرگ شود.
\end{itemize}

\menpage{۱۰}{pic10.jpg — 20261010\_053302370.jpg؛ شمارهٔ دست‌نویس در تصویر دیده نمی‌شود}
\section*{مننژیت؛ علائم بالینی}
\begin{itemize}
\item تب، سردرد، سفتی گردن، کاهش سطح هوشیاری؛ تهوع و استفراغ، فتوفوبی.
\item کرنیگ، برودزینسکی: نشانه‌های تحریک مننژ.
\item تریاد مننژیت: ۱. تب، ۲. سردرد، ۳. سفتی گردن.
\item سفتی گردن نشانهٔ پاتوگنومونیک تحریک مننژ است.
\item کاهش سطح هوشیاری از لتارژی تا کما: در ۷۵\% بیماران.
\item افزایش \en{ICP} از عوارض مننژیت است: بدتر شدن یا کاهش سطح هوشیاری، ادم پاپی، مردمک میدریاتیک با واکنش ضعیف به نور، فلج عصب ۶، وضعیت دسربره، رفلکس کوشینگ (۱. برادی‌کاردی، ۲. \en{$\uparrow$ BP}، ۳. تنفس نامنظم).
\item عامل اصلی آبتوندیشن و کما در بیمار مننژیت؛ مهم‌ترین عارضهٔ آن: هرنی مغزی.
\end{itemize}
\subsection*{آنالیز \en{CSF} در مننژیت باکتریال}
\begin{itemize}
\item فشار \en{CSF}: بیشتر از \en{$180\ \mathrm{mm\,H_2O}$}؛ گلوکز: \en{$<40\ \mathrm{mg/dl}$}.
\item \en{WBC}: ۱۰۰–۱۰٬۰۰۰ در هر \en{$\mu\mathrm{L}$} با ارجحیت نوتروفیل.
\item \en{RBC}: نباید باشد؛ \en{LP} تروماتیزه.
\item پروتئین: \en{$>45\ \mathrm{mg/dl}$}.
\item رنگ‌آمیزی گرم: در بیش از ۶۰\% موارد مثبت.
\item کشت: در بیش از ۸۰\% موارد مثبت.
\item \en{PCR}: کشف \en{DNA} باکتری.
\item نسبت قند \en{CSF} به سرم: نسبت \en{$<0.4$} به نفع مننژیت باکتریال حاد؛ سایر \en{DX}: مننژیت قارچی، سلی، کارسینوماتوز.
\end{itemize}

\menpage{۱۱}{pic11.jpg — 20261010\_053310987.jpg؛ شمارهٔ دست‌نویس: ۳}
\section*{مننژیت؛ برخورد و درمان تجربی}
\begin{itemize}
\item وقتی علائم بالینی به نفع مننژیت باکتریال است: سریعاً نمونهٔ کشت خون گرفته شود؛ شروع \en{AB} تجربی + دگزامتازون.
\item اورژانس است: \en{AB} باید ظرف ۶۰ دقیقه آغاز شود.
\item کدام بیماران قبل از \en{LP} نیاز به تصویربرداری دارند؟ ترومای اخیر به سر؛ نقص ایمنی؛ علائم فوکال نورولوژیک؛ ادم پاپی؛ کاهش سطح هوشیاری.
\end{itemize}
\subsection*{آنتی‌بیوتیک تجربی در مننژیت}
\begin{itemize}
\item شیرخوار پره‌ترم و نوزاد: آمپی‌سیلین + سفوتاکسیم.
\item شیرخوار ۱–۳ ماه: آمپی‌سیلین + سفتریاکسون یا سفوتاکسیم.
\item اطفال بالای ۳ ماه و بالغین زیر ۵۵ سال با ایمنی طبیعی: سفوتاکسیم یا سفتریاکسون یا سفپیم + ونکومایسین (\en{MRSA}).
\item بالای ۵۵ سال یا سوءمصرف الکل یا وجود بیماری‌های ناتوان‌کننده: آمپی‌سیلین + سفوتاکسیم یا سفتریا یا سفپیم (سودومونا) + ونکومایسین (\en{MRSA}).
\item عفونت بیمارستانی، پس از نوروسرجری، پس از تروما به سر، مبتلایان به \unclear{نوتروپنی}، اختلال ایمنی سلولی به هر علت: آمپی‌سیلین + سفتا یا مروپنم (سودومونا) + ونکومایسین (\en{MRSA}).
\item انسفالیت هرپسی در رأس \en{DDX} مننژیت باکتریال است: آسیکلوویر در همهٔ درمان‌های تجربی.
\item عوارض سفپیم: تشنج، میوکلونوس، انسفالوپاتی؛ اگر علائم وجود داشت سفپیم نمی‌دهیم.
\end{itemize}

\menpage{۱۲}{pic12.jpg — 20261010\_053320834.jpg؛ شمارهٔ دست‌نویس: ۴}
\section*{مننژیت؛ انتخاب آنتی‌بیوتیک}
\begin{itemize}
\item اگر پاتوژن احتمالی: استرپ پنومونیه، هموفیلوس آنفلوانزا، لیستریا مونوسایتوژن، باسیل‌های گرم منفی هوازی (سودومونا، \en{E. coli}): ونکو + مروپنم.
\item مروپنم را هیچ‌وقت به تنهایی نمی‌دهیم [حاشیه: سفتریا + ونکو].
\item شایع‌ترین عامل مننژیت در بیماران بدون طحال: استرپ پنومونیه.
\item مواردی که آمپی‌سیلین جهت پوشش لیستریا به درمان اضافه می‌شود: زیر ۳ ماه، بالای ۵۵ سال، شک به نقص ایمنی سلولی در اثر بیماری‌های مزمن، پیوند عضو، بارداری، بدخیمی، سرکوب ایمنی.
\item مواردی که مترونیدازول جهت پوشش بی‌هوازی و باسیل گرم منفی به درمان اضافه می‌شود: اوتیت، سینوزیت، ماستوئیدیت.
\end{itemize}
\subsection*{درمان مننژیت مننگوکوکی}
\begin{itemize}
\item کوکسی گرم منفی هوازی؛ پتشی و پورپورا.
\item نایسریا: گونه‌های حساس: پنی‌سیلین \en{G}؛ گونه‌های مقاوم: سفتریا یا سفوتاکسیم؛ طول مدت درمان ۱ هفته.
\end{itemize}
\subsection*{پروفیلاکسی مننژیت مننگوکوکی}
\begin{itemize}
\item در بیمار و تمام افرادی که تماس نزدیک با بیمار داشته‌اند.
\item ریفامپین: \en{$600\ \mathrm{mg}$ BD} برای ۲ روز؛ \en{$10\ \mathrm{mg/kg}$ BID} برای ۲ روز در کودک بالای ۱ سال.
\item تماس نزدیک: بوسیدن، اسباب‌بازی مشترک، نوشیدن و سیگار مشترک.
\item پروفیلاکسی حاملگی: ریفامپین ممنوع است؛ آزیترومایسین \en{$500\ \mathrm{mg}$} تک‌دوز؛ سفتریا \en{$250\ \mathrm{mg}$} تک‌دوز.
\end{itemize}

\menpage{۱۳}{pic13.jpg — 20261010\_053325513.jpg؛ شمارهٔ دست‌نویس: ۵}
\section*{درمان مننژیت پنوموکوکی}
\begin{itemize}
\item یک سفالوسپورین (سفتریا یا سفوتاکسیم یا سفپیم) + ونکومایسین.
\item طول درمان ۲ هفته وریدی.
\item تکرار \en{LP}، ۲۴–۳۶ ساعت پس از شروع درمان برای اطمینان از استریل شدن \en{CSF}.
\item [خط‌خورده:] \en{RF}: \unclear{پنومونی پنوموکوکی (مهم‌ترین)}.
\end{itemize}
\subsection*{درمان مننژیت لیستریایی}
\begin{itemize}
\item آمپی‌سیلین (اگر بیمار بدحال: + جنتامایسین).
\item طول درمان ۳ هفته.
\item آلرژی به پنی‌سیلین: \en{TMP-SMX} (کوتریموکسازول).
\item \en{RF} (مواردی که آمپی‌سیلین باید به درمان اضافه شود): زیر ۳ ماه، بالای ۵۵ سال، شک به نقص ایمنی سلولی در اثر بیماری مزمن، پیوند عضو، بارداری، بدخیمی، سرکوب ایمنی.
\item یک باسیل گرم مثبت هوازی متحرک است.
\item غذاهای آلوده: سالاد کلم، شیر، پنیر، غذاهای سریع آماده‌شده، سوسیس خوب پخته‌نشده، گوشت به‌میزان کافی پخته‌نشده.
\end{itemize}
\subsection*{دگزامتازون در مننژیت باکتریال}
\begin{itemize}
\item ۲۰ دقیقه قبل یا حداقل همزمان با شروع اولین دوز \en{AB}.
\item اگر بیش از ۶ ساعت از \en{AB} گذشته باشد دیگر اثر ندارد.
\item بیشترین اثر بر مننژیت پنوموکوکی.
\item کاهش عوارض نورولوژیک، کری حسی–عصبی، مرگ.
\end{itemize}

\menpage{۱۴}{pic14.jpg — 20261010\_053336054.jpg؛ شمارهٔ دست‌نویس: ۶}
\section*{ریسک‌فاکتورهای مرگ ناشی از مننژیت}
\begin{itemize}
\item سطح هوشیاری پایین در بدو ورود؛ تشنج ظرف ۲۴ ساعت اول؛ نشانه‌های \en{$\uparrow$ ICP}؛ شیرخوارگی و سن بالای ۵۰ سال؛ شوک یا نیاز به ونتیلاسیون؛ تأخیر در شروع درمان.
\item گلوکز \en{CSF}: \en{$<40\ \mathrm{mg/dl}$} و پروتئین \en{CSF}: \en{$>300\ \mathrm{mg/dl}$}: پیش‌آگهی بد.
\end{itemize}
\subsection*{مننژیت ویروسی حاد}
\begin{itemize}
\item سردرد، تب، علائم تحریک مننژ، نمای \en{CSF} التهابی.
\item سردرد معمولاً در ناحیهٔ فرونتال یا رترواربیت است، با فتوفوبی و درد در حرکات چشم همراهی دارد.
\item ریجید گردن خفیف است، فقط در انتهای \en{ROM} رخ می‌دهد.
\item بیماران اغلب هوشیار هستند؛ کاهش سطح هوشیاری بیشتر در انسفالیت ویروسی دیده می‌شود.
\item تشنج و نشانه‌های فوکال نورولوژیک وجود ندارد.
\end{itemize}
\subsection*{آنالیز \en{CSF}: مهم‌ترین تست در مننژیت ویروسی}
\begin{itemize}
\item پلئوسیتوز لنفوسیتی: ۲۵–۵۰۰ سلول در میکرولیتر.
\item پروتئین طبیعی یا اندکی بالا؛ گلوکز طبیعی.
\item \en{Opening P}: طبیعی یا کمی بالا (۱۰۰–۳۵۰).
\item ارگانیسمی یافت نمی‌شود.
\item مهم‌ترین تست تشخیص ویروس‌ها \en{PCR} است؛ در \en{HSV}، انتروویروس روش انتخابی است.
\item \en{PCR} مدفوع در انتروویروس کمک‌کننده است.
\end{itemize}

\menpage{۱۵}{pic15.jpg — 20261010\_053340605.jpg؛ شمارهٔ دست‌نویس: ۷}
\section*{مننژیت ویروسی}
\begin{itemize}
\item کشت ویروس حساسیت پایینی دارد.
\end{itemize}
\subsection*{انتروویروس‌ها}
\begin{itemize}
\item شایع‌ترین مننژیت‌های ویروسی؛ بیشتر تابستان و پاییز.
\item در کودکان زیر ۱۵ سال شایع‌تر است.
\item تب، سردرد، سفتی گردن، استفراغ، بی‌اشتهایی، اسهال، فارنژیت، میالژی.
\item \en{CSF}: پلئوسیتوز لنفوسیتی، گلوکز طبیعی، \en{pro} طبیعی یا اندکی بالا.
\item روش انتخابی تشخیص: \en{RT-PCR} بر روی \en{CSF}.
\item بیشترین حساسیت در ۴۸ ساعت اول؛ قابل انجام تا ۵ روز.
\item بر روی نمونهٔ گلو و مدفوع تا هفته‌ها مثبت.
\item درمان: حمایتی است.
\item در نوزادان و افراد هایپوگاما: عفونت مزمن و شدید.
\end{itemize}
\subsection*{هرپس سیمپلکس}
\begin{itemize}
\item دومین علت مننژیت ویروسی؛ در زنان شایع‌تر است.
\item عامل اصلی مننژیت \en{HSV-2}؛ عامل اصلی انسفالیت \en{HSV-1}.
\item \en{PCR} مایع \en{CSF} (\en{CSF} کلاسیک ویروس).
\item مننژیت \en{Mollaret}: مننژیت آسپتیک یا ویروسی راجعه؛ اکثراً علت \en{HSV}.
\item درمان هرپس: آسیکلوویر، فامسیکلوویر.
\end{itemize}

\menpage{۱۶}{pic16.jpg — 20261010\_053349895.jpg؛ شمارهٔ دست‌نویس: ۸}
\section*{مننژیت اوریونی}
\begin{itemize}
\item اواخر زمستان و اوایل بهار؛ در مردها شایع‌تر.
\item همراه با پاروتیت، اورکیت، اوفوریت، پانکراتیت، افزایش آمیلاز و لیپاز.
\item \en{CSF}: پلئوسیتوز با بیش از ۱۰۰۰ سلول؛ ۷۵\% لنفوسیت.
\item در ۱۰–۳۰\% موارد گلوکز کاهش می‌یابد (هایپوگلیکوراکی).
\end{itemize}
\subsection*{درمان مننژیت ویروسی}
\begin{itemize}
\item درمان علامتی: مسکن، تب‌بر، ضد استفراغ؛ اصلاح آب و الکترولیت‌ها.
\item اندیکاسیون بستری: نقص ایمنی، اختلال هوشیاری، تشنج، علائم فوکال نورولوژیک، \en{CSF} آتیپیک.
\end{itemize}
\subsection*{مننژیت تحت‌حاد}
\begin{itemize}
\item سل، قارچ (کریپتوکوکوس نئوفورمنس)، سیفلیس (شایع‌ترین).
\item اعصاب کرانیال ۷ و ۸ هستند.
\item علائم برای روزها و هفته‌ها: سردرد، سفتی گردن، تب \en{low grade}، \unclear{بی‌قراری}.
\item ممکن است اختلالات اعصاب کرانیال و تعریق شبانه هم داشته باشد.
\end{itemize}
\subsection*{مننژیت سلی}
\begin{itemize}
\item \en{CSF}: افزایش \en{Opening Pressure}؛ پلئوسیتوز لنفوسیتی (۱۰۰–۵۰۰)؛ افزایش میزان \en{pro} (۱–۵ \en{g/L})؛ کاهش گلوکز (۲۰–۴۰ \en{mg/dl}).
\item استاندارد طلایی تشخیص: کشت.
\item آخرین لولهٔ \en{LP} بهترین نمونه برای اسمیر \en{AFB} است.
\end{itemize}
\subsection*{مننژیت قارچی}
\begin{itemize}
\item \en{CSF}: پلئوسیتوز مونونوکلئر یا لنفوسیتی، \en{$\uparrow$ pro}، \en{$\downarrow$} گلوکز.
\item تشخیص: اسمیر \en{India ink}.
\end{itemize}

\menpage{۱۷}{pic17.jpg — 20261010\_053354881.jpg؛ شمارهٔ دست‌نویس: ۹}
\section*{مننژیت؛ تشخیص و درمان}
\begin{itemize}
\item در مننژیت \en{C. immitis}: ائوزینوفیل در \en{CSF}.
\item تست آنتی‌ژن پلی‌ساکارید کریپتوکوکی: بسیار حساس و اختصاصی.
\item مننژیت سیفلیسی: تشخیص: \en{FTA-ABS} سرم یا \en{MHA-TP} سرم مثبت بوده و در \en{CSF} پلئوسیتوز لنفوسیتی و \en{$\uparrow$ pro} دیده شود یا \en{VDRL} مایع \en{CSF} مثبت شود.
\end{itemize}
\subsection*{درمان}
\begin{itemize}
\item سل: ایزونیازید، ریفامپین، اتامبوتول، پیرازینامید.
\item در بیماران \en{HIV} مثبت مننژیت سلی: دگزا هم می‌دهیم.
\item قارچ: در بیمار \en{HIV} منفی و غیرپیوندی: آمفوتریسین \en{B} + فلوسیتوزین برای ۴ هفته؛ اگر عارضهٔ نورولوژیک: برای ۶ هفته.
\item بیماران پیوندی: آمفوتریسین لیپوزومال + فلوسیتوزین.
\item سیفلیس: پنی‌سیلین.
\end{itemize}
\ifnum\MeningitisStandalone=1 \expandafter\enddocument\fi
